% GENERATED FILE. Edit canonical lessons, then run npm run generate:books.
% canonical-source-hash: fnv1a64:9ac4218a9cb9b414
% canonical-lessons: ES-C602-voto, ES-C602-elecciones, ES-C602-candidato, ES-C602-votante, ES-C602-democracia

\chapter{Las elecciones}
\label{ch:es-las-elecciones}
\hlchaptermodality{602}

\begin{chapteropening}
\textbf{By the end of this chapter:} \emph{I can talk about an election, the candidates and the voters.}

Complete the last lesson of chapter 602: I can talk about an election, the candidates and the voters.
\end{chapteropening}

\section[el voto]{el voto — the vote}

\label{lesson:ES-C602-voto}



\begin{quote}
Before the new one: say the Spanish for the minister, then the Spanish for the town councillor.
\end{quote}

\subsection*{You'll want to know: el voto}
\textbf{el voto} — \textquotedblleft{}the vote\textquotedblright{}.

\begin{quote}
Cada voto es importante.
\end{quote}

\begin{quote}
\emph{Every vote matters.}
\end{quote}

\emph{Votar} — to vote.

\begin{cousinweb}
\emph{voto} is a learned borrowing from Latin \emph{votum} 'a vow, a wish', from \emph{vovere} 'to vow' (EtymDB).
\end{cousinweb}

\begin{grammarlens}[title={What you've built}]
One more word for telling, explaining and describing at B1.
\end{grammarlens}

\subsection*{Your turn}
\begin{itemize}
\raggedright
  \item \emph{Say it:} \emph{el voto}
  \item \emph{Say it:} \emph{el voto}, once more
  \item \emph{Say it:} a sentence of your own with \emph{el voto}
  \item \emph{Recall:} say the Spanish for the minister, then the Spanish for the town councillor, then say \emph{el voto} again
\end{itemize}

\begin{tcolorbox}[breakable,colback=teal!4,colframe=teal!35!black,title={Before you move on}]
What does el voto mean? (\textquotedblleft{}The vote\textquotedblright{}.) Say it once more.
\end{tcolorbox}
\section[las elecciones]{las elecciones — the elections}

\label{lesson:ES-C602-elecciones}



\begin{quote}
Before the new one: say the Spanish for the town councillor, then the Spanish for the vote.
\end{quote}

\subsection*{You'll want to know: las elecciones}
\textbf{las elecciones} — \textquotedblleft{}the elections\textquotedblright{}.

\begin{quote}
Las elecciones son en mayo.
\end{quote}

\begin{quote}
\emph{The elections are in May.}
\end{quote}

Plural for a political vote; in the singular \emph{la elección} is a choice.

\begin{cousinweb}
\emph{elecciones} is the plural of \emph{elección}, a learned borrowing from Latin \emph{electio} 'choice', built on \emph{eligere} 'to pick out', from \emph{legere} 'to gather' (EtymDB).
\end{cousinweb}

\begin{grammarlens}[title={What you've built}]
One more word for telling, explaining and describing at B1.
\end{grammarlens}

\subsection*{Your turn}
\begin{itemize}
\raggedright
  \item \emph{Say it:} \emph{las elecciones}
  \item \emph{Say it:} \emph{las elecciones}, once more
  \item \emph{Say it:} a sentence of your own with \emph{las elecciones}
  \item \emph{Recall:} say the Spanish for the town councillor, then the Spanish for the vote, then say \emph{las elecciones} again
\end{itemize}

\begin{tcolorbox}[breakable,colback=teal!4,colframe=teal!35!black,title={Before you move on}]
What does las elecciones mean? (\textquotedblleft{}The elections\textquotedblright{}.) Say it once more.
\end{tcolorbox}
\section[el candidato]{el candidato — the candidate}

\label{lesson:ES-C602-candidato}



\begin{quote}
Before the new one: say the Spanish for the vote, then the Spanish for the elections.
\end{quote}

\subsection*{You'll want to know: el candidato}
\textbf{el candidato} — \textquotedblleft{}the candidate\textquotedblright{}.

\begin{quote}
Hay tres candidatos para alcalde.
\end{quote}

\begin{quote}
\emph{There are three candidates for mayor.}
\end{quote}

A woman is \emph{la candidata}. Also used for a job.

\begin{cousinweb}
\emph{candidato} is a learned borrowing from Latin \emph{candidatus} 'clothed in white', built on \emph{candidus} 'white': Roman candidates wore white togas (Lewis \& Short).
\end{cousinweb}

\begin{grammarlens}[title={What you've built}]
One more word for telling, explaining and describing at B1.
\end{grammarlens}

\subsection*{Your turn}
\begin{itemize}
\raggedright
  \item \emph{Say it:} \emph{el candidato}
  \item \emph{Say it:} \emph{el candidato}, once more
  \item \emph{Say it:} a sentence of your own with \emph{el candidato}
  \item \emph{Recall:} say the Spanish for the vote, then the Spanish for the elections, then say \emph{el candidato} again
\end{itemize}

\begin{tcolorbox}[breakable,colback=teal!4,colframe=teal!35!black,title={Before you move on}]
What does el candidato mean? (\textquotedblleft{}The candidate\textquotedblright{}.) Say it once more.
\end{tcolorbox}
\section[el votante]{el votante — the voter}

\label{lesson:ES-C602-votante}



\begin{quote}
Before the new one: say the Spanish for the elections, then the Spanish for the candidate.
\end{quote}

\subsection*{You'll want to know: el votante}
\textbf{el votante} — \textquotedblleft{}the voter\textquotedblright{}.

\begin{quote}
Muchos votantes no fueron a votar.
\end{quote}

\begin{quote}
\emph{Many voters did not go to vote.}
\end{quote}

The same form for a woman: \emph{la votante}.

\begin{cousinweb}
\emph{votante} is built on \emph{votar}, which comes from \emph{voto}, a learned borrowing from Latin \emph{votum} 'a vow' (EtymDB).
\end{cousinweb}

\begin{grammarlens}[title={What you've built}]
One more word for telling, explaining and describing at B1.
\end{grammarlens}

\subsection*{Your turn}
\begin{itemize}
\raggedright
  \item \emph{Say it:} \emph{el votante}
  \item \emph{Say it:} \emph{el votante}, once more
  \item \emph{Say it:} a sentence of your own with \emph{el votante}
  \item \emph{Recall:} say the Spanish for the elections, then the Spanish for the candidate, then say \emph{el votante} again
\end{itemize}

\begin{tcolorbox}[breakable,colback=teal!4,colframe=teal!35!black,title={Before you move on}]
What does el votante mean? (\textquotedblleft{}The voter\textquotedblright{}.) Say it once more.
\end{tcolorbox}
\section[la democracia]{la democracia — the democracy}

\label{lesson:ES-C602-democracia}



\begin{quote}
Before the new one: say the Spanish for the candidate, then the Spanish for the voter.
\end{quote}

\subsection*{You'll want to know: la democracia}
\textbf{la democracia} — \textquotedblleft{}the democracy\textquotedblright{}.

\begin{quote}
La democracia necesita participación.
\end{quote}

\begin{quote}
\emph{Democracy needs participation.}
\end{quote}

\emph{Democrático} is the adjective.

\begin{cousinweb}
\emph{democracia} is a learned borrowing from Greek \emph{dēmokratia}, built from \emph{dēmos} 'the people' and \emph{kratos} 'power' (EtymDB).
\end{cousinweb}

\begin{grammarlens}[title={What you've built}]
One more word for telling, explaining and describing at B1.
\end{grammarlens}

\subsection*{Your turn}
\begin{itemize}
\raggedright
  \item \emph{Say it:} \emph{la democracia}
  \item \emph{Say it:} \emph{la democracia}, once more
  \item \emph{Say it:} a sentence of your own with \emph{la democracia}
  \item \emph{Recall:} say the Spanish for the candidate, then the Spanish for the voter, then say \emph{la democracia} again
\end{itemize}

\begin{tcolorbox}[breakable,colback=teal!4,colframe=teal!35!black,title={Before you move on}]
What does la democracia mean? (\textquotedblleft{}The democracy\textquotedblright{}.) Say it once more.
\end{tcolorbox}
